\documentclass[11pt]{article}

\usepackage[T1]{fontenc}
\usepackage{lmodern}
\usepackage{amsmath,amssymb,amsthm}
\usepackage[margin=1in]{geometry}
\usepackage[colorlinks=true,citecolor=blue,linkcolor=blue,urlcolor=blue]{hyperref}
\setlength{\emergencystretch}{2em}

\newtheorem{theorem}{Theorem}
\newtheorem{proposition}[theorem]{Proposition}
\newtheorem{lemma}[theorem]{Lemma}
\newtheorem{corollary}[theorem]{Corollary}

\newcommand{\Q}{Q}
\newcommand{\cA}{\mathcal A}
\newcommand{\cD}{\mathcal D}
\newcommand{\cU}{\mathcal U}

\title{A Componentwise Repair of the Two-Layer Lemma\\
in the Daykin--Frankl Argument}
\author{Xiaojian Zeng\\
\small PBC School of Finance, Tsinghua University, 100190\\
\small \texttt{xiaojianzeng@foxmail.com}}
\date{September 7, 2026}

\begin{document}
\maketitle

\begin{abstract}
Williams recently reduced the Daykin--Frankl width inequality for convex
subsets of a Boolean lattice to a two-layer lemma for an arbitrary finite
poset.  The posted proof of that lemma contains an undefined auxiliary set
and checks only one of two proposed antichains.  We prove the lemma by a
different construction: components of a bipartite comparison graph are
assigned between two output antichains, with isolated vertices handled by
the only possible boundary conflicts.  Both outputs are checked in every
comparison direction.  This establishes the lemma for all finite posets and
validates the stated induction, including empty sections.  We also give a
three-element-chain example showing that the natural set-theoretic
interpretation of the undefined auxiliary set does not repair the posted
sorting argument; the example does not contradict the lemma.
\end{abstract}

\section{Context and statement}

For a finite poset $S$, let $w(S)$ be the maximum size of an antichain in
$S$, with $w(\varnothing)=0$.  Products are ordered coordinatewise, and
$\Q_1=\{0<1\}$.  A subset of a poset is \emph{convex} if it contains every
point lying between two of its points.

Daykin and Frankl conjectured the Boolean-lattice width inequality now given
in Corollary~\ref{cor:DF} below \cite{DaykinFrankl1983}.  Duffus, Howard and
Leader proved it for initial segments of the binary order and developed
related width results for downsets \cite{DuffusHowardLeader2019}.  Williams's
recent preprint states the full inequality and reduces it to the following
two-layer assertion \cite[Lemma~2]{Williams2026}.  In the proof posted as
version~1, a set $C$ used in the sorting rules is not defined, and the second
of two output sets is declared an antichain by symmetry rather than checked
\cite[Section~3]{Williams2026}.  Our contribution is a complete replacement
for that sorting step.  We make no claim here about priority beyond this
specific repair.

Let $R$ be a finite poset and let $P\subseteq R\times\Q_1$ be convex.  Write
\[
  P=(P_0\times\{0\})\cup(P_1\times\{1\}),
  \qquad P_i=\{x\in R:(x,i)\in P\}.
\]

\begin{theorem}[Two-layer lemma]\label{thm:two-layer}
For every such $R$ and $P$,
\[
  2w(P)\geq w(P_0\times\Q_1)+w(P_1\times\Q_1).
\]
No nonemptiness, rankedness, or lattice assumption is required.
\end{theorem}

\section{The component construction}

Put $B=P_0\cap P_1$.  Convexity gives the elementary implication
\begin{equation}\label{eq:convexity}
 p\in P_0,\quad q\in P_1,\quad p\leq q
 \quad\Longrightarrow\quad p,q\in B.
\end{equation}
Indeed, $(p,1)$ and $(q,0)$ both lie between $(p,0)$ and $(q,1)$.

For $i\in\{0,1\}$, choose a maximum antichain of $P_i\times\Q_1$ and
write it uniquely by layers as
\[
 \cA_i=(L_i\times\{0\})\cup(H_i\times\{1\}).
\]
Thus $L_i$ and $H_i$ are antichains in $R$, and
\begin{equation}\label{eq:input-cross}
 \ell\in L_i,\ h\in H_i\quad\Longrightarrow\quad \ell\not\leq h.
\end{equation}

Form a bipartite graph $G$.  Its left vertices are occurrence copies
$L_1^*$ of the elements of $L_1$, its right vertices are occurrence copies
$H_0^*$ of the elements of $H_0$, and two occurrences are adjacent exactly
when their underlying values are comparable in $R$.  Occurrence copies keep
the two roles distinct when a value belongs to $L_1\cap H_0$.

\begin{lemma}[Uniform orientation]\label{lem:orientation}
Every nontrivial component of $G$ is exactly one of the following types:
\begin{enumerate}
\item a single equality edge $\ell=h$;
\item all its edges satisfy $\ell<h$ for $\ell\in L_1$ and $h\in H_0$;
\item all its edges satisfy $h<\ell$ for $\ell\in L_1$ and $h\in H_0$.
\end{enumerate}
\end{lemma}

\begin{proof}
An equality edge is an isolated component: any further neighbor would be
comparable with the equal value in one of the antichains $L_1$ or $H_0$.
Suppose next that an edge satisfies $\ell<h$.  Any other $L_1$-neighbor
$\ell'$ of $h$ must satisfy $\ell'<h$; the alternatives $h\leq\ell'$ and
$h=\ell'$ would make $\ell$ and $\ell'$ comparable.  Any other
$H_0$-neighbor $h'$ of $\ell'$ must then satisfy $\ell'<h'$; otherwise
$h'\leq\ell'<h$ would make $h'$ and $h$ comparable.  Propagating along
paths proves type~2.  Reversing the order proves type~3.
\end{proof}

We now assign every occurrence vertex to one of two sets of underlying
values, denoted $U$ and $D$.
\begin{itemize}
\item In a type~1 or type~2 component, assign all $L_1^*$ occurrences to
  $U$ and all $H_0^*$ occurrences to $D$ (the native assignment).
\item In a type~3 component, assign all $L_1^*$ occurrences to $D$ and all
  $H_0^*$ occurrences to $U$ (the swapped assignment).
\item Assign an isolated $\ell^*\in L_1^*$ to $D$ if some $a\in L_0$
  satisfies $a\leq\ell$, and otherwise to $U$.
\item Assign an isolated $h^*\in H_0^*$ to $U$ if some $b\in H_1$
  satisfies $h\leq b$, and otherwise to $D$.
\end{itemize}

Every layer move is legal.  In a type~3 component each vertex has an
incident edge $h<\ell$ with $h\in P_0$ and $\ell\in P_1$, so
\eqref{eq:convexity} puts both underlying values in $B$.  If an isolated
$\ell^*$ or $h^*$ is swapped, its defining boundary comparison and
\eqref{eq:convexity} again put its value in $B$.  Hence a swapped point
belongs to the required opposite section.

Define two subsets of $P$ by
\begin{equation}\label{eq:outputs}
 \cU=(L_0\times\{0\})\cup(U\times\{1\}),
 \qquad
 \cD=(D\times\{0\})\cup(H_1\times\{1\}).
\end{equation}

\begin{proof}[Proof of Theorem~\ref{thm:two-layer}]
We first check $\cU$.  The set $L_0$ is an antichain.  Within $U$, values
coming from $L_1^*$ are pairwise incomparable, as are values coming from
$H_0^*$.  A comparable cross-family pair would be an edge of $G$, but the
two sides of every component are assigned to opposite outputs.  Thus $U$ is
an antichain.

The only remaining possible comparison in $\cU$ has the form $a\leq u$
with $a\in L_0$ and $u\in U$.  If $u$ comes from a swapped $H_0^*$
occurrence, this is forbidden by \eqref{eq:input-cross} for $\cA_0$.  If it
comes from a native isolated $L_1^*$ occurrence, the comparison is forbidden
by the isolated-vertex rule.  If it comes from $L_1^*$ in a native equality
component, its $H_0^*$ neighbor has the same value; if it comes from a
type~2 component, it lies below an $H_0^*$ neighbor.  In either case
$a\leq u$ would imply $a\leq h$ for some $h\in H_0$, again contradicting
\eqref{eq:input-cross}.  These four cases exhaust $U$, so $\cU$ is an
antichain.

We check $\cD$ separately.  The graph argument just used shows that $D$ is
an antichain, and $H_1$ is an antichain.  A remaining comparison would have
the form $d\leq b$ with $d\in D$ and $b\in H_1$.  If $d$ comes from a
swapped $L_1^*$ occurrence, this is forbidden by
\eqref{eq:input-cross} for $\cA_1$.  If it comes from a native isolated
$H_0^*$ occurrence, it is forbidden by the isolated-vertex rule.  In a
native equality component, the value of $d$ equals that of its $L_1^*$
neighbor; in a type~2 component, an $L_1^*$ neighbor lies below $d$.
Thus either alternative would give $\ell\leq b$ for some $\ell\in L_1$,
contradicting \eqref{eq:input-cross}.  Hence $\cD$ is also an antichain.

Every occurrence in $L_1^*\sqcup H_0^*$ is assigned once.  No two
occurrences in the same output have the same underlying value: the only
possible cross-family equality is an equality edge, whose endpoints receive
opposite assignments.  Therefore
\[
 |U|+|D|=|L_1|+|H_0|.
\]
Using \eqref{eq:outputs},
\begin{align*}
 |\cU|+|\cD|
 &=|L_0|+|L_1|+|H_0|+|H_1|\\
 &=w(P_0\times\Q_1)+w(P_1\times\Q_1).
\end{align*}
Both outputs have size at most $w(P)$, proving the assertion.  The argument
also covers empty sections, isolated vertices, equality components, and a
disconnected or empty comparison graph.
\end{proof}

\section{The Boolean-lattice consequence}

Let $\Q_n=(2^{[n]},\subseteq)$, with $\Q_0$ the one-point poset.  The
standard width formula is
$w(\Q_m)=\binom{m}{\lfloor m/2\rfloor}$.

\begin{corollary}[Daykin--Frankl inequality]\label{cor:DF}
For all integers $n,k\geq0$ and every convex $P\subseteq\Q_n$,
\[
 w(P\times\Q_k)\geq w(\Q_{n+k})\,|P|\,2^{-n}.
\]
This includes $P=\varnothing$.
\end{corollary}

\begin{proof}
Induct on $n$ simultaneously for every $k\geq0$.  When $n=0$, either
$P=\varnothing$, giving zero on both sides, or $P=\Q_0$, giving
$w(\Q_k)=w(\Q_k)$.

Assume the assertion at $n$ for all $k$.  Identify
$\Q_{n+1}=\Q_n\times\Q_1$, let $P\subseteq\Q_{n+1}$ be convex, and write
its two sections as $P_0,P_1\subseteq\Q_n$.  Each section is convex.
Moreover, $P\times\Q_k$ is convex, and after permuting product coordinates
it is a convex subset of $\Q_{n+k}\times\Q_1$ with sections
$P_0\times\Q_k$ and $P_1\times\Q_k$.  Theorem~\ref{thm:two-layer} gives
\[
 w(P\times\Q_k)
 \geq\frac12\sum_{i=0}^1 w(P_i\times\Q_{k+1}).
\]
The induction hypothesis at the shifted parameter $k+1$ gives
\[
 w(P_i\times\Q_{k+1})
 \geq w(\Q_{n+k+1})|P_i|2^{-n}.
\]
Since $|P|=|P_0|+|P_1|$, averaging produces
$w(\Q_{n+k+1})|P|2^{-(n+1)}$, as required.
\end{proof}

\section{Why the natural undefined-set repair fails}

The following example concerns only the sorting construction in
\cite[Section~3]{Williams2026}, not Theorem~\ref{thm:two-layer}.

\begin{proposition}\label{prop:source-failure}
If the undefined set $C$ in that construction is interpreted as
\[
 C=(P_0\cup P_1)\setminus(P_0\cap P_1),
\]
then its first proposed output need not be an antichain, even when $R$ is a
three-element chain.
\end{proposition}

\begin{proof}
Let $R=\{0<1<2\}$, $P_0=\{0,1,2\}$, and $P_1=\{0,1\}$.  The resulting
$P$ is $R\times\Q_1$ with only $(2,1)$ removed; it is a downset and hence
convex.  Choose maximum input antichains
\[
 L_0=\{2\},\quad H_0=\{0\},\qquad
 L_1=\{1\},\quad H_1=\{0\}.
\]
Both section products have width two: the displayed sets give lower bounds,
and their two horizontal layers give chain covers of size two.

Here $B=P_0\cap P_1=\{0,1\}$ and $C=\{2\}$.  Under the posted rules, the
$L_1$ occurrence $1$ is sent to $L$ because no element of
$L_0\cup(H_0\cap C)$ lies below it.  The $H_0$ occurrence $0$ is also sent
to $L$, because the element $0\in H_1$ lies above it.  Consequently the
first proposed output contains
\[
 (2,0),\ (0,1),\ (1,1),
\]
and $(0,1)<(1,1)$.  Thus that output is not an antichain.  On this example
$w(P)=2$ and both section-product widths are two, so the two-layer inequality
itself holds with equality.
\end{proof}

\section{Exact finite audit and limitations}

As corroboration, an exact program enumerated every naturally labelled
finite poset through five elements, every convex subset of its product with
$\Q_1$, and all relevant widths by exact antichain dynamic programming.  It
tested 180,102 convex instances and found no counterexample.  Every
isomorphism class in this range occurs because every finite poset admits a
linear extension and can therefore be naturally labelled.  A second exact
program tested the component assignment for every pair of maximum input
antichains in those instances: all 1,684,785 pairs satisfied legality, both
antichain conditions, and cardinality preservation.  These finite checks do
not prove the universal theorem; the proof above does.

For methodological context only, the registered abstract of a separate
finite-poset enumeration preprint describes integral arithmetic and independent
evaluators for its stated endpoints \cite{Zeng2026SSRN7385118}.  That record
does not support either our universal theorem or the particular finite census
reported here; we cite it only as a nearby example of proof-oriented exact
computation on finite posets.

The argument establishes exactly the arbitrary-finite-poset lemma and the
displayed induction.  Proposition~\ref{prop:source-failure} invalidates only
one natural interpretation of the posted sorting rules.  We do not assert
that no other repair exists, nor do we make a general priority claim for the
Daykin--Frankl theorem.

\section*{Computational and AI assistance}

Language models assisted in discovering the comparison-graph construction,
developing exact verification code, and drafting the exposition.  The
universal proof was subsequently reconstructed and checked case by case in a
separate review context.  Exact programs independently enumerated the stated
finite ranges and reproduced the three-chain certificate.  No finite
computation is used as a substitute for the universal argument.

\bibliographystyle{plain}
\bibliography{references}

\end{document}
